\section{Large-prime labels and their short averages}\label{sec:labels}

We encode the large prime factors by two independently chosen multiplicative
labels.  Their one-variable distributions have limits independent of fixed
residue conditions.  We establish those limits and the weighted short-average
estimate needed to prove that the two labels decorrelate at consecutive
integers.  In the short-average estimate the original counting scale $x$ tends to
infinity before the auxiliary parameter $B$ does.
The scale $x>1$ is real unless a statement explicitly restricts it to integers.

\subsection{Labels and their one-variable distributions}

Fix an integer $J\geq2$.  For $1\leq k<J$ put
\begin{equation}\label{eq:labels-bins}
 \mathcal B_{k,x}=(x^{k/J},x^{(k+1)/J}].
\end{equation}
If $E$ is a set of primes, $\Omega_E(n)$ denotes the number of prime factors
of $n$ belonging to $E$, counted with multiplicity.  Fix complex numbers
$\zeta_1,\ldots,\zeta_{J-1}$ and $\xi_1,\ldots,\xi_{J-1}$ of absolute
value one, independent of $x$, with no relation required between the two
vectors.  Define the completely multiplicative function $f_x$ by
\begin{equation}\label{eq:labels-function}
 f_x(p)=
 \begin{cases}
  \zeta_k,&p\in\mathcal B_{k,x},\quad 1\leq k<J,\\
  1,&p\notin\bigcup_{1\leq k<J}\mathcal B_{k,x}.
 \end{cases}
 \qquad
 f_x(n)=\prod_{k=1}^{J-1}\zeta_k^{\Omega_{\mathcal B_{k,x}}(n)}.
\end{equation}
Define the completely multiplicative label associated with $(\xi_k)$ by
\begin{equation}\label{eq:labels-second-function}
 g_x(n)=\prod_{k=1}^{J-1}\xi_k^{\Omega_{\mathcal B_{k,x}}(n)}.
\end{equation}
Both labels have absolute value one.
For every fixed $C>0$, every integer $n\leq Cx$ has at most $J$ prime
factors in the union of the bins, once $x$ is sufficiently large in terms
of $C,J$.  Indeed, $J+1$ such factors would have product exceeding
$x^{1+1/J}>Cx$.

\begin{lemma}[One-variable distribution]\label{lem:labels-marginal}
For every fixed $0<\alpha<\beta<\infty$, positive integer $q$, and residue
class $a\pmod q$, the distribution of
\[
 \bigl(\Omega_{\mathcal B_{k,x}}(n)\bigr)_{1\leq k<J}
\]
among the integers $\alpha x<n\leq\beta x$, $n\equiv a\pmod q$,
normalized to have mass one, converges as $x\to\infty$.  Its limit depends
only on $J$, and not on $\alpha,\beta,q,a$.

Consequently there is a number $\mu=\mu(J,\zeta_1,\ldots,\zeta_{J-1})$
with $|\mu|\leq1$ such that
\begin{equation}\label{eq:labels-marginal-mean}
 \lim_{x\to\infty}
 \frac{q}{(\beta-\alpha)x}
 \sum_{\substack{\alpha x<n\leq\beta x\\n\equiv a\pmod q}}f_x(n)
 =\mu.
\end{equation}
The same assertion holds for $g_x$, with a mean $\mu_g$ depending only
on $J$ and the vector $(\xi_k)_{k=1}^{J-1}$.
Write
\begin{equation}\label{eq:labels-centered}
 F_x(n)=f_x(n)-\mu.
\end{equation}
Then $|F_x(n)|\leq2$, and for every fixed positive integer $u$,
\begin{equation}\label{eq:labels-multiplier}
 F_x(un)=F_x(n),\qquad g_x(un)=g_x(n)\qquad(n\geq1)
\end{equation}
for all sufficiently large $x$ in terms of $u,J$.
\end{lemma}

\begin{proof}
We prove convergence by computing all mixed factorial moments of the
bin counts.  We first record the elementary prime estimates used in this
calculation.  Put $\vartheta(z)=\sum_{p\le z}\log p$.  For $z\ge2$ we have
\begin{equation}\label{eq:distribution-prime-estimates}
 \vartheta(z)\ll z,\qquad
 \pi(z)\ll\frac z{\log z},\qquad
 \sum_{p\le z}\frac{\log p}{p}=\log z+O(1).
\end{equation}
Indeed, the primes in $(m,2m]$ divide $\binom{2m}{m}$, giving
$\vartheta(2m)-\vartheta(m)\ll m$.  Dyadic summation proves the first
estimate, and separating the primes at $\sqrt z$ proves the second.
Moreover,
\[
 \sum_{p^\nu\le z}\log p
 =\sum_{1\le\nu\le\log_2z}\vartheta(z^{1/\nu})\ll z,
\]
because the terms with $\nu\ge2$ are $O(\sqrt z\log z)$.
Expanding $\log n$ over prime powers now gives
\[
 \sum_{n\le z}\log n
 =\sum_{p^\nu\le z}\log p\,\lfloor z/p^\nu\rfloor
 =z\sum_{p\le z}\frac{\log p}{p}+O(z).
\]
The floor errors use the preceding prime-power bound, and the terms
with $\nu\ge2$ use
$\sum_p\sum_{\nu\ge2}(\log p)/p^\nu<\infty$.
Integral comparison on the left proves the third estimate in
\eqref{eq:distribution-prime-estimates}.  Partial summation yields
\begin{equation}\label{eq:distribution-prime-measure}
 \sum_{x^s<p\le x^t}\frac1p\longrightarrow\log(t/s)
 \qquad(0<s<t\ \text{fixed}),
\end{equation}
as well as $\sum_{p\le z}1/p=\log\log z+O(1)$.

The number of integers $n\leq\beta x$ divisible by $p^2$ for some
prime $p>x^{1/J}$ is at most
\[
 \beta x\sum_{p>x^{1/J}}\frac1{p^2}+O_\beta(\sqrt x)=o(x).
\]
This remains negligible relative to the size of any one fixed progression
in the lemma.  We may therefore replace multiplicity counts by counts of
distinct primes.

For the mixed factorial moments of these distinct-prime counts, write
$(b)_r=b(b-1)\cdots(b-r+1)$ and $(b)_0=1$.
Fix nonnegative integers $r_1,\ldots,r_{J-1}$, and let
$r=r_1+\cdots+r_{J-1}$.  The case $r=0$ is immediate.  For $r>0$, a term
in the factorial-moment expansion selects, in order, $r_k$ distinct
primes from bin $k$, for each $k$.  Denote their product by $d$; only
$d\leq\beta x$ can contribute.  The number of such selections is
$O_{J,\beta,r}(x/\log x)$.  To see this, fix every selected prime except
one, and write $d_1$ for their product.  If the last prime has any
admissible choice, its upper bound $\beta x/d_1$ exceeds $x^{1/J}$.
The prime-counting upper bound therefore gives at most
\[
 O_{J,\beta}\!\left(\frac{x}{d_1\log x}\right)
\]
choices.  Summing this bound over the other primes costs a bounded product
of reciprocal prime sums over the bins.  We may drop distinctness and
product restrictions when taking this upper bound.

For sufficiently large $x$, every selected prime is coprime to $q$.
The Chinese remainder theorem then gives
\[
 \#\{\alpha x<n\leq\beta x:n\equiv a\pmod q,\ d\mid n\}
 =\frac{(\beta-\alpha)x}{qd}+O(1).
\]
The sum of the errors over all selections is $o(x)$.  After normalization,
the factorial moment is consequently the reciprocal sum over the
selections with $d\leq\beta x$, up to $o(1)$.

In the variable $v=\log p/\log x$, \eqref{eq:distribution-prime-measure}
says that the reciprocal prime measure on bin $k$ converges to $dv/v$
on $(k/J,(k+1)/J]$.  The product restriction is
\[
 \sum_{\text{selected primes}}v\leq1+\frac{\log\beta}{\log x}.
\]
The limiting product measure is absolutely continuous, so the boundary
$\sum v=1$ has measure zero.  Repeated selections do not change the limit:
their reciprocal contribution is bounded by a constant times
$\sum_{p>x^{1/J}}p^{-2}=o(1)$.  Returning from distinct-prime counts to
multiplicity counts also costs $o(1)$ in the normalized moment, because
the exceptional set has size $o(x)$ and the counts are bounded.

List the selected bin indices as $k_1,\ldots,k_r$, with $r_k$
occurrences of $k$.  We have proved the explicit limit
\begin{equation}\label{eq:labels-factorial-limit}
\begin{aligned}
 &\lim_{x\to\infty}\frac{q}{(\beta-\alpha)x}
 \sum_{\substack{\alpha x<n\leq\beta x\\n\equiv a\pmod q}}
       \prod_{k=1}^{J-1}
          \bigl(\Omega_{\mathcal B_{k,x}}(n)\bigr)_{r_k}\\
 &\qquad =
 \int_{\substack{v_i\in(k_i/J,(k_i+1)/J]\ (1\le i\le r)\\
                  v_1+\cdots+v_r\le1}}
       \prod_{i=1}^r\frac{dv_i}{v_i}.
\end{aligned}
\end{equation}
For $r=0$, both sides are one.  The right-hand side depends only on
$J$ and the multi-index $(r_k)$.  All count vectors lie in a fixed
finite set, and mixed factorial polynomials span the functions on that
set.  The asserted distributional convergence follows.

Taking the expectation of the function
$(b_k)\mapsto\prod_k\zeta_k^{b_k}$ gives \eqref{eq:labels-marginal-mean}; using the
vector $(\xi_k)$ gives the mean $\mu_g$.  Finally, if every prime factor
of $u$ is at most $x^{1/J}$, then $f_x(u)=g_x(u)=1$.
Complete multiplicativity of $f_x$ and $g_x$ proves
\eqref{eq:labels-multiplier}.
\end{proof}

The centered function $F_x$ is not asserted to be multiplicative.
Its eventual invariance under each fixed multiplier is the exact property
in \eqref{eq:labels-multiplier} that will be used below.

The same distributional limit holds for averages over $1\leq n\leq x$.
Indeed, apply the lemma to any fixed function of the count vector on
$\alpha x<n\leq x$, then let $\alpha\downarrow0$; that function is
bounded on the finite set of possible vectors, so the omitted interval
has normalized contribution $O(\alpha)$.  In particular,
the limit in \eqref{eq:labels-factorial-limit} is unchanged when its
normalized progression sum is replaced by
$x^{-1}\sum_{1\le n\le x}\prod_k(\Omega_{\mathcal B_{k,x}}(n))_{r_k}$.
This gives \eqref{eq:labels-marginal-mean} on the initial segment as
well.  A fixed shift of the integer argument changes only a bounded
number of terms after nonpositive arguments are omitted, so the label
means have the same limits after such a shift.

\subsection{The mixed correlation and its short-average input}

The joint law will follow by finite Fourier inversion from the following
statement for the two independently chosen phase vectors.  The factor
$F_x(n+1)$ is centered, while $g_x(n)$ has absolute value one.

\begin{proposition}[Mixed decorrelation]\label{prop:mixed}
For every fixed integer $J\geq2$ and every pair of fixed phase vectors
$(\zeta_k)_{k=1}^{J-1}$ and $(\xi_k)_{k=1}^{J-1}$ of absolute value one,
the corresponding labels satisfy
\begin{equation}\label{eq:mixed}
 \lim_{x\to\infty}\frac1x\sum_{1\leq n<x}
                   \overline{g_x(n)}F_x(n+1)=0,
\end{equation}
where $F_x=f_x-\mu$ and the limit is through all integer scales.
\end{proposition}

After establishing the short-average input below, we will suppose that
the average in \eqref{eq:mixed} stays a fixed positive distance from zero
along a sequence and derive a contradiction.  The input is used in
Section~\ref{sec:conclusion} at the final step of the proof.

For each auxiliary parameter $B$, let $P_B$ be a finite set of primes,
and suppose that $\min P_B\to\infty$ as $B\to\infty$.  An empty set
may be allowed by interpreting its minimum as $+\infty$ and its
product as one.  For a fixed real
$v\geq0$, define
\begin{equation}\label{eq:labels-G}
 G_{B,v}(m)=\prod_{\substack{p\mid m\\p\in P_B}}
                   \left(\frac12+\frac12p^{-v/B}\right).
\end{equation}
This is a real, nonnegative, 1-bounded multiplicative function.  For
fixed $B$ it is periodic, with period dividing $\prod_{p\in P_B}p$.
To see its fair-split meaning, put
$S_B(m)=\{p\in P_B:p\mid m\}$ and form a random product $b$ by
including each prime of $S_B(m)$ independently with probability $1/2$.
The empty product is one.  Independence gives
\[
 G_{B,v}(m)
 =\mathbb E_{\mathrm{split}}\!\left[\exp\!\left(-\frac{v\log b}{B}\right)\right].
\]
These are the fair-split averages used to approximate the endpoint
features in Section~\ref{sec:conclusion}.

\begin{lemma}[Weighted short averages of the centered labels]
\label{lem:labels-short}
Let $T=T(B)\to\infty$, let $L_B$ be positive integers tending to
infinity, and let $a_B$ be any integer for each $B$.  Fix
$0<s_0<s_1<\infty$, a positive integer $D$, a residue $r\pmod D$,
and $v\geq0$.  Then
\begin{equation}\label{eq:labels-short}
 \lim_{B\to\infty}\limsup_{x\to\infty}
 \frac1{Tx}\sum_{s_0Tx\leq n\leq s_1Tx}
 \left|
   \frac1{L_B}\sum_{i=a_B+1}^{a_B+L_B}
     \ind_{i\equiv r\ (D)}F_x(n+i)G_{B,v}(n+i)
 \right|^2=0.
\end{equation}
For each fixed $B$ all the arguments of the arithmetic functions are
positive once $x$ is sufficiently large.  The same conclusion holds
if the inner upper limit is taken along any sequence $x\to\infty$.
\end{lemma}

The proof interpolates the centered function of the finitely many bin
counts by real, nonnegative multiplicative functions.  Once the residue of
the origin $n$ is fixed, the congruence on $i$ becomes a fixed residue
condition on the argument $m=n+i$.  Resolving that condition by Dirichlet
characters produces two different tasks.
For the principal character, short averages must be compared with long
means whose centered linear combination vanishes.  For a nonprincipal
character, the short averages themselves must be small.  We record the two
published inputs for these tasks and prove the required distance estimate
before returning to the weighted lemma.

\subsection{The short-interval inputs}

For a bounded arithmetic function $g$ and $H>0$, set
\[
 \mathcal A_Hg(z)=\frac1H\sum_{z<n\leq z+H}g(n),
 \qquad
 \mathcal L_Xg=\frac1X\sum_{X<n\leq2X}g(n).
\]
For a 1-bounded multiplicative function $g$, define
\begin{equation}\label{eq:labels-distance}
 \mathbb D(g,n^{it};X)^2
 =\sum_{p\leq X}\frac{1-\operatorname{Re}(g(p)p^{-it})}{p},
 \qquad
 M(g;X)=\inf_{|t|\leq X}\mathbb D(g,n^{it};X)^2.
\end{equation}

We use the following two forms of the short-interval theorems.  Their
uniformity in the multiplicative function is part of the statements.

\begin{theorem}[Real short-interval comparison]
\label{thm:labels-real-MR}
There is a function $r(H)$ tending to zero as $H\to\infty$ such that,
for every real multiplicative function $g:\N\to[-1,1]$ and
$2\leq H\leq X$,
\begin{equation}\label{eq:labels-real-MR}
 \frac1X\int_X^{2X}
   |\mathcal A_Hg(z)-\mathcal L_Xg|^2\,dz\leq r(H).
\end{equation}
\end{theorem}

This is a consequence of the uniform exceptional-set estimate in
\cite[Theorem~1]{MatomakiRadziwill2016}, using boundedness to pass to
mean square.  The terms in that estimate depending on $X$ can be
absorbed into $r(H)$ because $X\geq H$.

\begin{theorem}[Complex short averages]
\label{thm:labels-complex-MRT}
For every multiplicative function $g:\N\to\C$ with $|g|\leq1$ and
$X\geq H\geq10$,
\begin{equation}\label{eq:labels-complex-quantitative}
 \frac1X\int_X^{2X}|\mathcal A_Hg(z)|^2\,dz
 \ll (1+M(g;X))e^{-M(g;X)}
       +\left(\frac{\log\log H}{\log H}\right)^2
       +(\log X)^{-1/50},
\end{equation}
with an absolute implied constant.
\end{theorem}

This is \cite[Theorem~A.1]{MatomakiRadziwillTao2015}, in the revised
version with the corrected proof of Proposition~A.3.  That proof permits
the factor $e^{-M}$; we use the weaker $(1+M)e^{-M}$, which also covers
$M<1$.  Thus divergence of the minimum over $|t|\leq X$ suffices for
our application.

The preceding statements also hold with integer origins, at the cost of
an error tending to zero with $H$.  Indeed, for $|g|\leq1$,
\[
 |\mathcal A_Hg(z)-\mathcal A_Hg(\lfloor z\rfloor)|\leq\frac2H.
\]
Integrating over unit intervals proves the claim, with $O(1/X)$ errors
at the endpoints.  Changing a short interval endpoint by a bounded
amount similarly costs $O(1/H)$.

\subsection{A uniform estimate for the character twists}

\begin{lemma}[Fixed nonprincipal characters]
\label{lem:labels-character-distance}
Fix $J\geq2$, positive constants $c,C$, a nonprincipal Dirichlet
character $\chi$ of modulus $q_\chi$, and a finite set $E$ of primes.  Suppose that, for
each $x$, $h_x$ is a 1-bounded multiplicative function satisfying
\[
 h_x(p)=\chi(p)\qquad(p\leq x^{1/J},\ p\notin E).
\]
With $X=cx$, one has
\begin{equation}\label{eq:labels-character-distance}
 \inf_{|t|\leq CX}\mathbb D(h_x,n^{it};X)^2\longrightarrow\infty
 \qquad(x\to\infty).
\end{equation}
The assertion is uniform over all the functions $h_x$ with this prime
agreement.  The threshold for $x$ may depend on $J,c,C,\chi,E$.
\end{lemma}

\begin{proof}
Put $y=x^{1/J}$ and $\sigma=1+1/\log x$.  For large $x$ we have
$y<X$.  Since the summands defining the distance are nonnegative,
we can restrict to $p\leq y$ and delete $E$ and the primes dividing
the modulus of $\chi$.  Deleting those finitely many primes costs only
a constant in the lower bounds below.

We record a prime-sum comparison, uniform in any factors of absolute
value at most one.  Replacing the weight $1/p$ on $p\leq y$ by
$p^{-\sigma}$ on all primes has absolute error $O_J(1)$.  Below $y$,
the error is at most
\[
 \frac1{\log x}\sum_{p\leq y}\frac{\log p}{p}=O_J(1).
\]
For the tail, partial summation and the prime-counting bound in
\eqref{eq:distribution-prime-estimates} give
\[
 \sum_{p>y}p^{-\sigma}
 \ll\int_y^\infty\frac{u^{-1-1/\log x}}{\log u}\,du
 =\int_{1/J}^\infty\frac{e^{-v}}v\,dv=O_J(1).
\]
The prime-power terms in the logarithm of either a zeta or a Dirichlet
$L$ Euler product are also uniformly $O(1)$ for $\sigma>1$.

For $|t|\leq1$ it follows that
\[
 \operatorname{Re}\sum_{p\leq y}\frac{\chi(p)p^{-it}}p
 =\log|L(\sigma+it,\chi)|+O_J(1)\leq C_{J,\chi}.
\]
Here we used only an upper bound for the logarithm: a nonprincipal
Dirichlet $L$-function has no pole at $1$ and is bounded on the compact
region under consideration.  Since \eqref{eq:distribution-prime-measure}
and the accompanying partial-summation estimate give
$\sum_{p\leq y}1/p=\log\log x+O_J(1)$, this proves a lower bound
$\log\log x-O_{J,\chi,E}(1)$ for this frequency range.

For $|t|>1$, let $m$ be the order of $\chi$ on the units.  The elementary
inequality $|1-z^m|\leq m|1-z|$, for $|z|=1$, gives
\[
 1-\operatorname{Re}(\chi(p)p^{-it})
 \geq\frac1{m^2}\bigl(1-\operatorname{Re}(p^{-imt})\bigr)
 \qquad(p\nmid q_\chi).
\]
The Vinogradov--Korobov bound, with any fixed logarithmic exponent
strictly between $2/3$ and $1$, gives
\begin{equation}\label{eq:labels-zeta-bound}
 |\zeta(\sigma+iu)|\ll
       (\log(|u|+3))^{3/4}
 \qquad(\sigma\geq1,\ |u|\geq1).
\end{equation}
The bound on the line $1$ follows from
\cite[(1.2)]{FordZeta2022}.  Its extension to the right can be seen
by applying the Phragm\'en--Lindel\"of principle in $1\leq\sigma\leq2$
to
\[
 \frac{s-1}{s+1}\,
       \frac{\zeta(s)}{(\log(s+3))^{3/4}}.
\]
The pole at $1$ is removed, the logarithmic power has an analytic
branch in this strip, and both vertical boundaries are bounded.
For $|\operatorname{Im}s|\geq1$, the removed rational factor is
bounded away from zero, which gives \eqref{eq:labels-zeta-bound}.
For $\sigma\geq2$ the same estimate follows from absolute
convergence.
Applying the prime-sum comparison just proved, we obtain uniformly for
$1<|t|\leq CX$,
\begin{align*}
 \sum_{p\leq y}\frac{1-\operatorname{Re}(p^{-imt})}{p}
 &=\log\log x-\log|\zeta(\sigma+imt)|+O_J(1)\\
 &\geq\tfrac14\log\log x-O_{J,c,C,\chi}(1).
\end{align*}
Combining the two frequency ranges proves
\[
 \mathbb D(h_x,n^{it};X)^2
 \geq\frac1{4m^2}\log\log x-O_{J,c,C,\chi,E}(1)
 \qquad(|t|\leq CX),
\]
which implies \eqref{eq:labels-character-distance}.
\end{proof}

\subsection{Proof of the weighted short-average lemma}

\begin{proof}[Proof of Lemma~\ref{lem:labels-short}]
Ter\"av\"ainen~\cite[Section~4, proof of Theorem~1.11]{Teravainen2018}
uses real multiplicative generating functions for large-prime counts and
recovers event coefficients from them.  Here the finite count range gives
a pointwise tensor interpolation of the centered function of all bin
counts by completely multiplicative functions.  We prove this
interpolation with coefficients independent of $B,x$.
Choose $J+1$ distinct numbers in $(0,1]$.  The Vandermonde matrix whose
entries are their powers of orders $0,\ldots,J$ is invertible.
Taking tensor products over the $J-1$ count coordinates shows that
there are finitely many complex constants $c_\ell$ and vectors
$\boldsymbol z_\ell=(z_{\ell,1},\ldots,z_{\ell,J-1})\in(0,1]^{J-1}$
such that
\begin{equation}\label{eq:labels-interpolation}
 \prod_{k=1}^{J-1}\zeta_k^{b_k}-\mu
 =\sum_\ell c_\ell\prod_{k=1}^{J-1}z_{\ell,k}^{b_k}
 \qquad(0\leq b_k\leq J).
\end{equation}
The coefficients and vectors depend only on the fixed labels and $J$.
Define
\[
 h_{\ell,x}(m)=\prod_{k=1}^{J-1}
                   z_{\ell,k}^{\Omega_{\mathcal B_{k,x}}(m)}.
\]
Each $h_{\ell,x}$ is real, nonnegative, 1-bounded, and completely
multiplicative.  On every interval $m\leq C_Bx$ with $C_B$ fixed
for fixed $B$, the identity
$F_x(m)=\sum_\ell c_\ell h_{\ell,x}(m)$ holds for all sufficiently
large $x$.

We next prove the needed estimate for each fixed residue class of the
integer argument.  For a class $r'\pmod D$, put
$d=\gcd(r',D)$ and $q=D/d$, and write $m=dk$.  Then $m\equiv r'\pmod D$
is equivalent to $k\equiv r'/d\pmod q$, a unit class modulo $q$.
For all sufficiently large $B$, no prime factor of $d$ belongs to
$P_B$, so $G_{B,v}(dk)=G_{B,v}(k)$.  For fixed $B$ and sufficiently
large $x$, \eqref{eq:labels-multiplier} also gives
$F_x(dk)=F_x(k)$, and the same equality holds for every interpolant
$h_{\ell,x}$.

By character orthogonality, the restricted $k$-sum is a fixed linear
combination of sums of
\begin{equation}\label{eq:labels-twisted-interpolants}
 h_{\ell,x}(k)G_{B,v}(k)\chi(k),\qquad \chi\pmod q.
\end{equation}
All these functions are 1-bounded and multiplicative.  The sum
length after division by $d$ is $H_B=L_B/d$.  Normalization by
$L_B$ is $d^{-1}$ times normalization by $H_B$.
Changing integer endpoints introduces an error $O_D(L_B^{-1})$.

Consider first the principal character $\chi_0\pmod q$.
The functions in \eqref{eq:labels-twisted-interpolants} are then real,
so Theorem~\ref{thm:labels-real-MR} applies to each interpolant.
Although their individual long means need not vanish, their centered
linear combination does.  More precisely, on any dyadic interval
$[X,2X]$ with $X=c_Bx$ and $c_B>0$ fixed for fixed $B$,
\[
 \sum_\ell c_\ell\mathcal L_X
       (h_{\ell,x}G_{B,v}\chi_0)
 =\frac1X\sum_{X<k\leq2X}F_x(k)G_{B,v}(k)\chi_0(k)
 \longrightarrow0.
\]
To justify the last limit, split the sum into residue classes modulo
the fixed common period of $G_{B,v}$ and $\chi_0$, and apply
Lemma~\ref{lem:labels-marginal} to every class.  The number and the
sizes of these residue classes may depend on $B$; they are fixed in
this $x$-limit.  The mean square of the centered short average thus
has inner upper limit at most a fixed multiple of $r(H_B)$, where
the multiple depends on the interpolation coefficients and $D$, but not on
$B$.

For a nonprincipal $\chi$, the function in
\eqref{eq:labels-twisted-interpolants} equals $\chi(p)$ at all
primes $p\leq x^{1/J}$ outside the finite set $P_B$.
Lemma~\ref{lem:labels-character-distance} therefore applies for
each fixed $B$, with $X=c_Bx$ and the frequency range $|t|\leq X$.
For all sufficiently large $B$ we have $H_B\geq10$, so
Theorem~\ref{thm:labels-complex-MRT} applies once $x$ is sufficiently
large.  Taking the inner upper limit in
\eqref{eq:labels-complex-quantitative} leaves at most a constant times
$(\log\log H_B/\log H_B)^2$.  This tends to zero as $B\to\infty$.
There are only finitely many interpolants and characters, so the same
holds after their linear combination.

Finally, for a given origin $n$, the condition $i\equiv r\pmod D$
in \eqref{eq:labels-short} is the condition
$m=n+i\equiv n+r\pmod D$.  There are only $D$ possible argument
residues $r'$, and it suffices to sum the bounds just proved for
these possibilities.  After $m=dk$, the short-interval origin is
$(n+a_B)/d$.  For fixed $B$ and all sufficiently large $x$,
\[
 \frac{s_0T}{2d}x\le\frac{n+a_B}{d}\le\frac{2s_1T}{d}x.
\]
Choose a dyadic cover of the fixed scaled interval
$[s_0T/(2d),2s_1T/d]$.  It gives intervals
$[c_{B,j}x,2c_{B,j}x]$ with every $c_{B,j}>0$ fixed in the inner
$x$-limit, and with a number of intervals bounded in terms of
$s_0,s_1,D$ independently of $B$.  Their scales are comparable to
$Tx/d$.  Mapping the origins to their integer parts has multiplicity
at most $d+1$, and replacing an origin by its integer part costs
$O(H_B^{-1})$.  The integer-origin version
of the preceding estimates consequently applies.  Normalizing by
$Tx$ instead of the length of each dyadic range changes only fixed
factors.  Since $H_B=L_B/d\to\infty$ for every $d\mid D$, this
proves \eqref{eq:labels-short}.

All assertions used an unrestricted upper limit as $x\to\infty$;
restricting that upper limit to a subsequence preserves them.
\end{proof}

\subsection{A profile of a nonzero mixed correlation}

We begin the proof by supposing that Proposition~\ref{prop:mixed} fails.
Fix an offending $J$ and pair of phase vectors.  Boundedness permits
passing to a subsequence on which the mixed correlation has a nonzero
limit:
\begin{equation}\label{eq:labels-bad-correlation}
 \lim_{x\to\infty}
  \frac1x\sum_{1\leq n<x}
                    \overline{g_x(n)}F_x(n+1)\ne0.
\end{equation}
The arithmetic amplification in the following sections will contradict
this limit; Section~\ref{sec:conclusion} completes the argument.

Let $\whZ$ be the profinite integers, with Haar probability
measure $dw$, and identify every integer with its natural image in
$\whZ$.  On $(0,\infty)\times\whZ$ define the locally finite measures
\[
 \lambda_x=\frac1x\sum_{n\ge1}\delta_{(n/x,n)},\qquad
 \nu_x=\frac1x\sum_{n\ge1}
       \overline{g_x(n)}F_x(n+1)\delta_{(n/x,n)}.
\]
Unweighted scale counting gives $\lambda_x\to dt\,dw$ against compactly
supported continuous tests.  This follows first for a continuous test in $t$ times
a residue-class indicator in $w$ by progression counting, and then
for every compactly supported continuous test by uniform approximation.
Also $|\nu_x|\le2\lambda_x$.  The resulting uniform variation bounds
on compact sets, weak compactness on a countable exhaustion, and a
diagonal extraction give a further subsequence along which $\nu_x$ converges
against these tests to a locally finite complex measure $\nu$.

For every compactly supported continuous $\Phi$, the weak convergences give
\[
 \left|\int\Phi\,d\nu\right|
 =\lim_{x\to\infty}\left|\int\Phi\,d\nu_x\right|
 \le2\lim_{x\to\infty}\int|\Phi|\,d\lambda_x
 =2\int_0^\infty\int_{\whZ}|\Phi(t,w)|\,dw\,dt.
\]
The dual characterization of variation implies $|\nu|\le2\,dt\,dw$.
The Radon--Nikodym theorem therefore gives a measurable function
\begin{equation}\label{eq:labels-profile}
 W:(0,\infty)\times\whZ\longrightarrow\C,
 \qquad |W(t,w)|\leq2\quad\text{almost everywhere},
\end{equation}
such that
\begin{equation}\label{eq:labels-profile-convergence}
 \lim_{x\to\infty}\frac1x\sum_{n\geq1}
       \overline{g_x(n)}F_x(n+1)\Phi(n/x,n)
 =\int_0^\infty\int_{\whZ}\Phi(t,w)W(t,w)\,dw\,dt
\end{equation}
for every compactly supported continuous $\Phi$ on the product.
Here and henceforth limits in $x$ may use the fixed subsequence.

Boundedness also permits replacing the compactly supported
continuous tests in \eqref{eq:labels-profile-convergence} by
the indicator of $0<t<1$: truncate near $0$, approximate the
remaining interval at its endpoints, and use the uniform bound
on both weighted counting measures and $W$.  Equation
\eqref{eq:labels-bad-correlation} therefore implies
\[
 \int_0^1\int_{\whZ}W(t,w)\,dw\,dt\ne0.
\]
Approximating this interval indicator by nonnegative smooth
functions supported in $(0,1)$, we can fix a real nonnegative
$\phi\in C_c^\infty((0,\infty))$ with
\begin{equation}\label{eq:labels-beta}
 \beta_*=
 \left|\int_0^\infty\int_{\whZ}
             \phi(t)W(t,w)\,dw\,dt\right|>0.
\end{equation}
The rest of the proof derives a contradiction from this fixed
profile and bump.  Every auxiliary choice will be made after
$J,f_x,g_x,F_x,W,\phi$ and the subsequence have been fixed.
